\chapter{Conclusão}

Os resultados obtidos ao longo desta prática demonstraram uma excelente concordância entre os valores teóricos e experimentais, validando a modelagem matemática através da função de transferência $H(s)$ para o filtro de realimentação múltipla (MFB). A análise em regime permanente senoidal permitiu prever com precisão o comportamento do circuito frente a variações de frequência.

A caracterização do filtro evidenciou sua natureza passa-faixa, com a frequência de ressonância experimental situada em 480,2~Hz, apresentando um desvio mínimo em relação ao valor teórico projetado[cite: 58, 67]. Embora o ganho máximo experimental ($|H| \approx 42,43$) tenha sido ligeiramente inferior ao valor teórico de $49,97$, a resposta global do sistema seguiu fielmente as curvas de magnitude e fase calculadas anteriormente. Tais discrepâncias são atribuídas às tolerâncias dos componentes passivos, especialmente dos capacitores, e às limitações práticas do amplificador operacional real.

A prática cumpriu plenamente seus objetivos, proporcionando:
\begin{itemize}
	\item Consolidação do cálculo analítico e determinação da função de transferência para circuitos ativos complexos;
	\item Domínio da análise de resposta em frequência e regime permanente senoidal;
	\item Experiência no uso de ferramentas de computação simbólica, como o WxMaxima, para a resolução de sistemas de equações lineares;
	\item Aperfeiçoamento no uso do osciloscópio e gerador de sinais para medições precisas de amplitude e defasagem;
	\item Compreensão da influência dos polos e zeros do sistema na estabilidade e filtragem de sinais.
\end{itemize}

Por fim, destaca-se que o estudo de filtros ativos e funções de transferência é fundamental para o projeto de sistemas de processamento de sinais e condicionamento de dados em engenharia. A validação experimental realizada através das medições sistemáticas no intervalo de 40~Hz a 11~kHz constitui uma base sólida para o desenvolvimento de sistemas eletrônicos mais sofisticados e robustos.